\documentclass[11pt,a4paper]{ctexart}

\usepackage{amsmath,amssymb,amsfonts}
\usepackage{geometry}
\usepackage{booktabs}
\usepackage{array}
\usepackage{enumitem}
\usepackage{hyperref}
\usepackage{float}
\usepackage{siunitx}
\usepackage{multirow}

\geometry{left=2.5cm,right=2.5cm,top=2.5cm,bottom=2.5cm}

\hypersetup{
	colorlinks=true,
	linkcolor=blue,
	urlcolor=blue,
	citecolor=blue
}

\title{阶段性成果综述：基于SSEPC模糊成对约束的电力用户聚类改进}

\begin{document}
	
	\maketitle
	
	
	\tableofcontents
	\newpage
	
	\section{加项式FCM变体与电力用户聚类的结合}
	
	\subsection{加项式半监督变体}
	
	本项目复刻 Luo et al.（2019, IEEE Trans. Smart Grid）的电力零售套餐推荐系统原论文将FCM作为黑箱直接引用（仅引Bezdek 1984，无任何参数细节。
	
	标准FCM的目标函数为：
	\begin{equation}
		J = \sum_{i}\sum_{j} u_{ij}^{m}\|x_{i} - c_{j}\|^{2}
		\quad (\text{式1, s.t. 行和为1, } u \ge 0)
	\end{equation}
	其中 $u_{ij}$ 为样本 $i$ 对簇 $j$ 的隶属度，$c_j$ 为簇中心，$m$ 为模糊指数，$x_i$是第i个样本的特征向量，$c_j$是第j个聚类中心的特征向量。在FCM家族谱系中，“往目标函数加辅助信息项”形成了一整条正则化/半监督分支（Pedrycz加项式、FCM\_S空间邻域项、SSFPC成对约束项等），统称“加项式”变体：
	\[
	J = \text{FCM项} + \beta \cdot \text{辅助信息项}.
	\]
	本项目主线采用的SSFPC（Wang et al. 2022, IEEE TFS）即该家族中“隶属度空间（$u$空间）成对约束”分支的代表。
	
	\subsection{与电力用户聚类的结合}
	
	数据来源于 Mendeley Data，DOI: 10.17632/zm4f727vvr.1 。原论文（Luo et al. 2019, IEEE TSG）使用SGSC项目 的 1,000 户居民数据（2013-01 至 2014-01，其参考文献 [30] 为 SGSC 官网）。本数据集是同一项目数据的清洗版。
	
		FCM项和辅助信息项分别由不同的数据产生，具体如下：
	
	\begin{enumerate}[label=(\arabic*)]
		\item FCM项: 就是标准的FCM目标函数，其中样本的特征向量包含了五个维度的信息（月均用电量、夏冬季月均、春秋季月均、工作时间周均、 非工作时间周均），从数据集 30 分钟原始负荷表格提取。
		\item 辅助信息项：区别于原论文SSFPC的监督者人工标注两样本相似程度$s_{pq}$ ，本文从社会画像自动推导 $s_{pq}$，6,031户全样本可得、零人工标注。数据集包含了样本用户的社会信息，如收入、燃气用量、房型、气候区。由于这些数据是混合类型，使用Gower算法来计算它们的相似度。这样辅助信息项与FCM项由不同的数据产生，避免了自己监督自己的情况，也符合半监督里用两两用户相似度来约束的原则。
		\item 看起来契合“社会信息过滤的主题”：本文的聚类模式为画像相近$\Rightarrow$ 用电行为相近 $\Rightarrow$ 隶属度分布相近，与论文标题“社会信息过滤（Social Information Filtering）”的语义一致；
		\item 场景与验证：落在电力零售套餐推荐场景，可能会为后续的推荐环节提供便利。
	\end{enumerate}
	
	\section{SSFPC的数学原理与迭代算法}
	
	\subsection{目标函数的构造与最小化}
	
	SSFPC在原版FCM目标函数的基础上叠加模糊成对约束项（2），综合考虑复杂度和聚类效果，暂且取模糊指数 $\gamma=2$、$\alpha=0$：
	\begin{align}
		\min_{\{u_i,c_j\}} \quad & \sum_{i=1}^{m}\sum_{j=1}^{k} \left(u_{ij}^{\gamma} - \alpha\right)\|x_i - c_j\|^{2}
		+ \beta \sum_{(p,q)\in N_s} C(u_p,u_q) \\
		\text{s.t.} \quad & \sum_{j=1}^{k} u_{ij} = 1, \quad \forall i=1,\ldots,m, \\
		& u_{ij} \ge 0, \quad \forall i=1,\ldots,m,\ \forall j=1,\ldots,k.
	\end{align}
	其中约束代价函数 $C$ 按 $s_{pq}$ 的符号采用两种不同的二次型构造：
	\begin{equation}
		C(u_p,u_q) =
		\begin{cases}
			\dfrac{1}{2}s_{pq}\|u_p - u_q\|^{2}, & s_{pq} \ge 0,\\[6pt]
			-s_{pq}\,u_p^{\top}u_q, & \text{otherwise}.
		\end{cases}
	\end{equation}
	
	符号约定：$\mathcal{S}$ 为模糊成对约束集合，每个元素记录一对样本间的相似程度；$\mathbb{N}_{\mathcal{S}}$ 为 $\mathcal{S}$ 的索引集；$s_{pq}\in[-1,1]$ 为样本 $p,q$ 间的相似度（块内样本间的约束强度下文简记为 $s_{ij}$）；$m$ 为样本总数；$k$ 为聚类数；$u_{ij}$ 为样本 $i$ 对簇 $j$ 的隶属度，$u_i = (u_{i1}, \ldots, u_{ik})^{\top}$ 为其隶属度向量；$u_p, u_q$ 特指被某条约束关联的一对样本的隶属度向量；$x_i$ 为样本 $i$ 的特征向量；$c_j$ 为簇 $j$ 的聚类中心。
	
	\(s_{pq}\) 是监督者对样本对 \((x_p, x_q)\) 相似性或相异性的模糊判断。其符号决定了引导方向：\(s_{pq}>0\) 对应相似约束（must-link 的模糊推广），要求把两者的隶属度向量拉近；\(s_{pq}<0\) 对应相异约束（cannot-link 的模糊推广），要求把两者推开。其绝对值 \(|s_{pq}|\) 决定了引导强度：\(|s_{pq}|\) 越大，两个对象在空间上被拉得越近或推得越远。分段函数的设计正是为了在执行这一引导方向与强度的同时，不破坏模糊特性。
	
	\(\frac{1}{2}s_{pq}\|u_p-u_q\|^{2}\) 负责把两者的隶属度向量拉近。在线性规划求解过程中，最小化 \(\|u_p-u_q\|\) 是要求两个隶属度向量（\(u_p, u_q\)）尽可能接近，\textbf{并不要求它们必须完全相等，更不要求它们退化为 one-hot 向量}。因此，两个模糊向量可以在保持模糊分布（例如 \((0.6,0.3,0.1)^\top\)）的同时彼此靠近，模糊特性得以保留。
	
	用最小化内积 \(u_p^\top u_q\) 来实现相异，当两个模糊向量在\textbf{不同分量上都有较大值}时，内积较小，满足“相异”；但每个向量本身仍可以是模糊的（例如 \(u_p=(0.7,0.2,0.1)^\top\)，\(u_q=(0.1,0.2,0.7)^\top\)），不会被迫退化为 one-hot。负约束用 \(-s_{pq}u_p^{\top}u_q\) 把两者推开（因 \(s_{pq}<0\)，该项系数为正，最小化它等价于压小内积，内积越小越相异）。当约束全部非负时（本文正约束特例），负约束分支不出现，目标退化为：
	\begin{equation}
		J = \sum_i \sum_j u_{ij}^{2}\|x_i - c_j\|^{2}
		+ \frac{\beta}{2}\sum_{(p,q)\in N_s} s_{pq}\|u_p - u_q\|^{2}.
	\end{equation}
	
	
	\subsection{交替迭代（MEM）}
	
	
	为简化计算，本文取 $\alpha = 0$：原论文的中心更新以 $\sum_i (u_{ij}^{2} - \alpha)$ 为分母，若其不大于零则判定该簇为空并删除；取 $\alpha = 0$ 时分母即 $\sum_i u_{ij}^{2}$，由隶属度行和为 $1$ 知其恒大于零，从机制上排除了删簇分支。求解采用交替优化（原论文称 MEM），每轮迭代只用两套更新：
	
	\begin{enumerate}[label=(\arabic*)]
		\item \textbf{中心更新（闭式）}。约束项不含簇中心，故中心更新不受加项影响，仍为隶属度加权均值：
		\begin{equation}\label{eq:c-update}
			c_j = \frac{\sum_i u_{ij}^{2} x_i}{\sum_i u_{ij}^{2}}.
		\end{equation}
		
		\item \textbf{隶属度更新（分为有闭式解和无闭式解）}。按样本是否受约束（即某一个样本是否和其他样本有相似度，如果和所有其他样本的相似度都为0，就认为这个样本和其他样本是解耦的，是不受约束的）分为两个子问题：
		
		不受约束样本（$i \notin \mathbb{N}_{\mathcal{S}}$）的隶属度仍有闭式解，即标准 FCM 更新：
		\begin{equation}\label{eq:u-closed}
			u_i = \frac{1}{\sum_{j=1}^{k} \|x_i - c_j\|^{2}}
			\left( \frac{1}{\|x_i - c_1\|^{2}}, \ldots, \frac{1}{\|x_i - c_k\|^{2}} \right)^{\!\top}.
		\end{equation}
		相对的，受约束样本的隶属度子问题无闭式解。为了简化计算，我们可以按照某些标准将集合 $\mathcal{S}$ 划分为若干个互不相交的子集 $\mathcal{S}_1, \mathcal{S}_2, \ldots, \mathcal{S}_h$，其中每个索引集中的样本彼此直接或间接关联。由于目标关于 $\{u_i\}$ 可分，该子问题随之分解为 $h$ 个相互独立的块子问题，记块 $\mathcal{S}_l$ 所含样本数 $r = |\mathbb{N}_{\mathcal{S}_l}|$：
		\begin{equation}\label{eq:block-qp}
			\begin{aligned}
				\min_{\{u_i \mid i \in \mathbb{N}_{\mathcal{S}_l}\}}\ &
				\sum_{i \in \mathbb{N}_{\mathcal{S}_l}} \sum_{j=1}^{k} u_{ij}^{2} \|x_i - c_j\|^{2}
				+ \frac{\beta}{2} \sum_{0 < s_{pq} \in \mathcal{S}_l} s_{pq} \left( \sum_{j=1}^{k} u_{pj}^{2} + \sum_{j=1}^{k} u_{qj}^{2} \right) \\
				& - \beta \sum_{0 > s_{pq} \in \mathcal{S}_l} s_{pq}\, u_p^{\top} u_q \\
				\text{s.t.}\ & \sum_{j=1}^{k} u_{ij} = 1,\quad \forall i \in \mathbb{N}_{\mathcal{S}_l}, \\
				& u_{ij} \ge 0,\quad \forall i \in \mathbb{N}_{\mathcal{S}_l},\ \forall j = 1, \ldots, k.
			\end{aligned}
		\end{equation}
		
	对每个块，把块内 \(r\) 个样本的隶属度向量按样本堆叠成一个长向量
	\[
	\mathbf{u}=(u_1^{\top},\ldots,u_r^{\top})^{\top}\in\mathbb{R}^{rk},
	\]
	于是该块的目标函数（距离项与约束项）可以合并写成一个二次型 $\frac{1}{2}\mathbf{u}^{\top}D\mathbf{u}$，同时每个样本的隶属度之和为 $1$ 这一要求可以写成矩阵形式 $A\mathbf{u}=\mathbf{1}^r$，其中 $D\in\mathbb{R}^{rk\times rk}$ 为对称矩阵，其分块构造见第~\ref{subsec:solver}~节。这样，每个块子问题就化归为一个\textbf{带行和与非负约束的二次规划}。该问题按黑塞矩阵 $D$ 的定性分两档求解：全部约束非负时 $D$ 恒正定，问题塌缩为单次凸二次规划（CQPP）；一般情形下 $D$ 可能不定，则采用对角块坐标下降（DBCD）逐样本求解——原论文证明了 DBCD 收敛到该问题的一个驻点，$D$的矩阵化构造、定性判定的逻辑链与上述两档求解的细节见第~\ref{subsec:solver}~节。
	\end{enumerate}
	
	\begin{table}[H]
		\centering
		\small
		\begin{tabular}{p{13.5cm}}
			\toprule
			\textbf{算法 1（MEM 外层循环）}：交替优化关于 $\{u_i\}$ 与 $\{c_j\}$ 的两组变量 \\
			\midrule
			\textbf{输入}：特征数据 $\{x_i\}$、模糊成对约束 $\mathcal{S}$（$s_{pq}$ 由 demographic 画像自动生成）、聚类数 $k$、参数 $\beta>0$、容差 $tol$。\\
			\textbf{输出}：隶属度 $\{u_i\}$ 与聚类中心 $\{c_j\}$。\\
			1：初始化 $u^{(0)}$（各样本取均匀隶属度 $(1/k,\ldots,1/k)^{\top}$ 或随机可行值），置 $t=1$。\\
			2：\textbf{中心更新}：固定 $u^{(t-1)}$，按式~\eqref{eq:c-update} 计算 $c_j^{(t)}$（$j=1,\ldots,k$）。\\
			3：\textbf{隶属度更新}：固定 $\{c_j^{(t)}\}$；不受约束样本按式~\eqref{eq:u-closed} 更新；每个约束块按式~\eqref{eq:block-qp} 求解——块QP凸时用 CQPP 求解器，不定时用 DBCD（内部再循环至其自身收敛）。\\
			4：若 $\|u^{(t)} - u^{(t-1)}\| < tol$，停止并输出 $u^{(t)}$；否则置 $t \leftarrow t+1$，回到第 2 步。\\
			\bottomrule
		\end{tabular}
	\end{table}
	
	\subsection{块二次规划的定性与求解}\label{subsec:solver}
	
	\begin{enumerate}[label=(\arabic*)]
		\item \textbf{正约束特例（v2版）}：$s_{pq}\ge 0$ 时块QP的黑塞矩阵恒严格正定（有定理与每块 $\lambda_{\min}$ 数值验证双保险），原论文的求解流程塌缩为单次凸QP，用OSQP（稀疏ADMM + 热启动）求解。
		\item \textbf{完整版含负约束（v3版）}：沿用式~\eqref{eq:block-qp} 的块划分，将块内 \(r\) 个样本的隶属度向量按样本堆叠后，块子问题可重写为标准二次规划形式：
		\begin{equation}\label{eq:qp-matrix}
			\min_{\mathbf{u}}\ \frac{1}{2}\,\mathbf{u}^{\top} D\, \mathbf{u}
			\quad \text{s.t.}\quad A\mathbf{u} = \mathbf{1}^{r},\quad \mathbf{u} \ge \mathbf{0}^{rk},
		\end{equation}
		其中 \(A \in \mathbb{R}^{r \times rk}\) 为行和约束的系数矩阵（每行对应一个样本的 \(\sum_j u_{ij} = 1\)），\(D \in \mathbb{R}^{rk \times rk}\) 为对称矩阵，按 \(k\) 行 \(k\) 列分块：对角块为
		\begin{equation}\label{eq:D-diag}
			D_{ii} = \mathrm{diag}\!\left( \|x_i - c_1\|^{2}, \ldots, \|x_i - c_k\|^{2} \right) 
			+ \frac{\beta}{2} \sum_{s_{ij} > 0} s_{ij} I,
		\end{equation}
		非对角块为
		\begin{equation}\label{eq:D-offdiag}
			D_{ij} = -\frac{\beta}{2}\,\mathrm{diag}(s_{ij}, \ldots, s_{ij}), \quad i \ne j,
		\end{equation}
		其中 \(I\) 为 \(k\) 阶单位阵，\(s_{ij}\) 为块内样本 \(i,j\) 间的约束强度。正、负约束同时进入同一个二次型 \(D\)：正约束经对角块中的 \(s_{ij}I\) 项体现，负约束经非对角块中的内积项体现；正是负约束使 \(D\) 可能不定，从而使式~\eqref{eq:qp-matrix} 不再一定是凸二次规划。
		
		对实对称矩阵，正定、半正定、不定完全由特征值符号决定：全部特征值为正（非负、既有正又有负）分别对应正定（半正定、不定），且正定等价于二次型 \(\mathbf{u}^{\top}D\mathbf{u}\) 对任意非零 \(\mathbf{u}\) 恒正。需注意"对角元全正"仅保证迹为正，并不足以保证正定——特征值的符号同时受非对角元影响。具体到 \(D\)：其迹由对角块（距离平方 \(+\) 正约束强度）构成，恒大于零，故负定被排除，只剩半正定或不定两种可能；而其不定性恰来自负约束——当 \(s_{ij}<0\) 时非对角块 \(D_{ij}=-\frac{\beta}{2}\mathrm{diag}(s_{ij})\) 为正，过强的正耦合会产生负特征值（最简单的例子：\(D=\bigl(\begin{smallmatrix}1&2\\2&1\end{smallmatrix}\bigr)\)，特征值为 \(3\) 与 \(-1\)）。反之，若约束全部非负，则非对角块非正，\(D\) 恒半正定——这正是 v2 正约束特例中块QP恒凸的来源。实务上按最小特征值 \(\lambda_{\min}\) 的符号判定定性：\(\lambda_{\min}>0\) 正定，\(\lambda_{\min}<0\) 不定；本文对每块 \(D\) 逐一作 \(\lambda_{\min}\) 数值验证。据此，若 \(D\) 半正定，式~\eqref{eq:qp-matrix} 即 CQPP，可直接得到其全局解；否则式~\eqref{eq:qp-matrix} 为 IQPP，一般的 IQPP 算法要么形式不匹配、要么计算量过大，但利用其特殊结构仍可有效求解——采用对角块坐标下降（DBCD）直接求解。
		
		具体地，DBCD 是一种坐标下降：把 $rk$ 维的联合优化拆成 $r$ 个 $k$ 维子问题，每次只更新一个样本的隶属度。扫描按 Gauss--Seidel 方式进行——按样本顺序 $1,\ldots,r$ 依次更新，更新第 $i$ 个样本时，排在前面的样本采用本轮刚算出的新值、排在后面的样本仍用上一轮的旧值（更新就地生效），一整轮扫过全部样本后进入下一轮，如此反复直至收敛。每次固定其余 \(r-1\) 个样本、仅解当前样本的 \(k\) 维子问题：
		\begin{equation}\label{eq:qp-sample}
			\begin{aligned}
				\min_{u_i}\ & \frac{1}{2}\,u_i^{\top} D_{ii}\, u_i
				+ (u_1^{\top}, \ldots, u_{i-1}^{\top}, u_{i+1}^{\top}, \ldots, u_r^{\top})
				(D_{i1}, \ldots, D_{i(i-1)}, D_{i(i+1)}, \ldots, D_{ir})^{\top} u_i \\
				\text{s.t.}\ & \mathbf{1}^{\top} u_i = 1, \quad u_i \ge \mathbf{0}^{k}.
			\end{aligned}
		\end{equation}
		该子问题恒为凸二次规划（CQPP），其原因在于不定性与样本间耦合绑定。把完整二次型按样本展开：
		\[
		\frac{1}{2}\mathbf{u}^{\top}D\mathbf{u}
		= \sum_{i}\frac{1}{2}u_i^{\top}D_{ii}u_i + \sum_{i<j}u_i^{\top}D_{ij}u_j,
		\]
		固定其余样本后，所有交叉项 \(u_i^{\top}D_{ij}u_j\)（\(j\neq i\)）退化为 \(u_i\) 的线性项，二次项只剩下 \(\frac{1}{2}u_i^{\top}D_{ii}u_i\)。而 \(D\) 的不定性恰来自样本间耦合（负约束经非对角块产生的正耦合撞出负特征值），\(D_{ii}\) 由构造恒正定（距离平方与正约束强度之和），且 \(D\) 的负特征值方向必然同时涉及多个样本、在单样本方向上不存在负曲率。因此子问题~\eqref{eq:qp-sample} 的黑塞矩阵恒正定，其全局解可由 CQPP 求解器获得；由于每个子问题都精确解到全局最小，目标值随扫描单调不增，DBCD 循环直至相邻两轮解的变化小于容差（通常取 \(10^{-3}\)），并收敛到式~\eqref{eq:qp-matrix} 的一个驻点（原论文证明了该收敛性）。
		
		在实现上，我们采用逐样本 Gauss--Seidel 串行扫描；每个样本的 \(k\) 维子问题以精确水填充算法解出，即按线性项排序枚举活跃前缀并以二分法兜底。该实现严格对齐官方 MATLAB code\footnote{\url{https://github.com/gamer1882/FDC}} 中的 \texttt{QPP\_IDsolver}/\texttt{Solver\_x0} 结构。
	\end{enumerate}
	
	\textbf{偏差声明}：官方代码私加了 $m/s$ 归一化因子（论文式18/19中并无此项），本文不采用；$\beta$ 一律按论文式(5)(6)口径报告。完整推导链与验证记录见项目内《SSFPC构造说明.md》及各版本README。
	
	\section{数据集来源与预处理}
	
	\subsection{数据来源}
	
	主数据为澳大利亚SGSC（Smart Grid Smart City）项目清洗版数据集（Roberts et al. 2019, Mendeley Data, DOI: 10.17632/zm4f727vvr.1），含6,031户新南威尔士州居民2013全年完整的30分钟级用电负荷（17,520个采样点无缺失）及配套人口统计（demographic）问卷字段。该数据与原论文所用1,000户同属SGSC项目但更大（6倍规模）；实验对比一律以本数据上自跑基线为准，不直接引用原论文数值。
	
	\subsection{预处理管线一：负荷 $\rightarrow$ 用电特征}
	
	按原论文Table I从30分钟负荷序列提取5维特征：月均用电量、夏季/冬季（南半球12--2月/6--8月）月均、春秋季月均、工作时间（澳大利亚惯例工作日9:00--17:00）周均、非工作时间周均，再经z-score标准化、按demographic行序对齐，供聚类输入（输出 \texttt{data/output/features\_std.npy}）。
	
	\subsection{预处理管线二：社会画像 $\rightarrow$ 两两相似度约束}
	
	从demographic表选取36个画像字段（有序型如收入/燃气档位、名义型如房型/气候区/地理分层、二值型如HAS-\*设施拥有情况、普查码与计量配置计数等），用Gower相似度处理混合类型并保留有序渐进性；映射为 $s = 2S - 1 \in [-1,1]$ 后稀疏化：每户取Top-10相似邻居构成正约束（42,927对，21个连通块），Bottom-10构成负约束（59,760对，$s\in[-0.871,0)$）。字段纪律：排除源自负荷曲线的字段（与5维特征冗余）与试验分组字段（非画像属性）。正+负边合并后全图为1个连通块。
	
	\section{A/B/C/D四组对照实验结果}
	
	\subsection{实验设计}
	
	四组配置均在6,031户 $\times$ 5维标准化特征上运行，$k=5$，$\beta=0.1$（A组 $\beta=0$ 即纯FCM），每组3个随机种子；指标分四层：模糊内部指标（PC/PE/XB）、硬划分内部指标（silhouette/Calinski--Harabasz/Davies--Bouldin）、约束满足度（正对同簇率/lift/残差）、外部参照（与Roberts预聚类标签的ARI/NMI，仅背景参照）与稳定性（seed间ARI）。
	
	\begin{table}[H]
		\centering
		\small
		\begin{tabular}{p{2.2cm}p{2.8cm}p{2.6cm}p{4.2cm}}
			\toprule
			组别 & 约束 & 求解器 & 定位 \\
			\midrule
			A\_fcm & 无（$\beta=0$） & 闭式更新 & 纯FCM基线 \\
			B\_gower15 & 旧版15字段正约束（54,473对） & OSQP凸QP & 约束有效性对照 \\
			C\_gower36 & 新版36字段正约束（42,927对） & OSQP凸QP & 字段扩充效果 \\
			D\_gower36neg & C + 负约束（59,760对） & 串行DBCD & 负约束消融 \\
			\bottomrule
		\end{tabular}
	\end{table}
	
	\subsection{主指标结果（3 seeds均值）}
	
	\begin{table}[H]
		\centering
		\small
		\setlength{\tabcolsep}{4pt}
		\begin{tabular}{lccccc}
			\toprule
			指标 & A纯FCM & B旧约束 & C新约束 & D加负约束 \\
			\midrule
			PC $\uparrow$ & 0.689 & 0.493 & 0.507 & 0.504 \\
			PE $\downarrow$ & 0.601 & 0.935 & 0.909 & 0.914 \\
			XB $\downarrow$ & 0.216 & 0.332 & 0.298 & 0.392 \\
			Silhouette $\uparrow$ & 0.449 & 0.423 & 0.433 & 0.395 \\
			CH $\uparrow$ & 9810 & 9345 & 9602 & 8854 \\
			DB $\downarrow$ & 0.810 & 0.833 & 0.821 & 0.856 \\
			正对同簇率 & 0.376 & 0.411 & 0.420 & 0.415 \\
			Lift（对随机基线）$\uparrow$ & 1.320 & 1.459 & 1.475 & 1.541 \\
			约束残差 $\downarrow$ & 0.620 & 0.227 & 0.270 & 0.266 \\
			目标函数终值 & 2392.14 & 3352.01 & 3240.51 & 3791.87 \\
			迭代轮数/单轮耗时量级 & $\approx$214轮/秒级 & $\approx$178轮/约102\,s & $\approx$178轮/约194\,s & 200（上限）/约261\,s \\
			\bottomrule
		\end{tabular}
	\end{table}
	
	注：B与C的约束集合不同，残差绝对值不可跨组直接比较，应看lift与同簇率；D组200轮为迭代上限到点，但目标尾部变化已在fp32噪音量级（相对约 $1\times10^{-7}$），视为数值收敛，且三个seed的objective\_ascents全为0。
	
	\subsection{负约束专项（C $\leftarrow$ D消融）}
	
	\begin{table}[H]
		\centering
		\small
		\begin{tabular}{lccc}
			\toprule
			指标 & C组（仅正约束） & D组（正+负） & 随机基线 \\
			\midrule
			负对同簇率 $\downarrow$ & 0.305 & 0.165 & 0.269 \\
			负对隶属度内积 $\downarrow$ & 0.300 & 0.161 & — \\
			负约束残差 $\downarrow$ & 0.136 & 0.069 & — \\
			簇规模（seed 0） & [1926,105,591,2189,1220] & [1792,133,705,2077,1324] & — \\
			\bottomrule
		\end{tabular}
	\end{table}
	
	\subsection{结论}
	
	\textbf{结论一（A$\rightarrow$B/C）：约束有用。} 加入社会画像约束后，正对同簇率lift从1.32提升到1.48，约束残差从0.62降至0.27，隶属度分布确实向“画像相似者隶属度分布相近”的方向移动；PC/PE变化反映划分更模糊，属约束项起作用的预期现象。
	
	\textbf{结论二（B$\rightarrow$C）：字段扩充全面更优。} 36字段版在七项可比指标上全胜15字段版，且相似度区分度问题（“大家都像”）得到缓解——$s$唯一值达79,196个，满分平局对占比从 $3.1\%$ 降至 $0.21\%$。
	
	\textbf{结论三（C$\rightarrow$D）：负约束有效且代价小。} cannot-link对子同簇率 $30.5\% \rightarrow 16.5\%$（远低于随机基线 $26.9\%$），隶属度内积减半，“推迟”真实发生；正约束满足度反而更好（lift $1.475\rightarrow 1.541$）；内部指标仅轻微下降（silhouette $-0.04$）。三组实验seed间ARI均为1.0（D组 $1.0/0.9565/0.9565$），结果稳定可复现。
	
	\section{结论与下一步工作}
	
	聚类改进线已全部打通：约束生成（Gower画像相似度，含正/负两档）、正约束特例凸QP求解、完整版串行DBCD求解、四层指标评估体系与A/B/C/D对照实验，三条核心结论均获数据支持。
	
	下一步（按优先级）：
	
	\begin{enumerate}[label=(\arabic*)]
		\item $k/\beta$网格扫描与敏感性分析（$\beta$需按约束密度重定标，$k$的取值依据待从原论文核实）；
		\item 推荐层实现——评分推断（5/3/2/1/0规则）、相似度加权聚合、Top-$N$输出与MAE验证，其中套餐评分数据来源（合成/powertochoose采集/按原论文4条选择模式仿真）待定，是当前最大未知数；
		\item 写法B（特征空间标量注入）消融对照。
	\end{enumerate}
	
\end{document}